%*******************************************************
% Declaration
%*******************************************************
\pdfbookmark[0]{Déclaration}{declaration}
\chapter*{Déclaration}
\label{ch:declaration}
\thispagestyle{empty}

\wip{Déclaration}{Vérifier d'abord ce que l'USMB et l'école doctorale exigent. Déclaration d'originalité usuelle; travaux en collaboration ou en encadrement identifiés comme tels; et la mention d'usage d'outils génératifs, à accorder mot pour mot avec le colophon.}

% ============================ TODO texte (~10 lignes) ==============================
% Le contenu classicthesis d'origine (« Put your declaration here. ») a ete
% retire: il etait encore en place.
%
%  1. VERIFIER CE QUE L'USMB ET L'ECOLE DOCTORALE EXIGENT reellement dans cette
%     page. La formulation ci-dessous n'a de sens que si rien n'est impose.
%
%  2. DECLARATION D'ORIGINALITE usuelle: travaux presentes, part personnelle,
%     travaux realises en collaboration ou en encadrement clairement identifies
%     comme tels.
%
%  3. ASSISTANCE PAR OUTILS GENERATIFS. De plus en plus d'etablissements
%     francais attendent une mention explicite. Si c'est le cas ici, une phrase
%     factuelle suffit: de tels outils ont ete employes pour la mise en forme
%     et la redaction de certaines parties de l'etat des lieux
%     (\cref{pt:context}) a partir de vos propres notes et des sources citees,
%     ainsi que pour la production du code de simulation; la responsabilite du
%     contenu, des demonstrations et des resultats reste entiere.
%     Declaree, la chose est banale; non declaree puis devinee, elle ne l'est
%     plus.
%
%  4. A APPARIER avec la passe de couverture bibliographique sur les blocs
%     numerotes de la \cref{pt:context}: les deux repondent a la meme question
%     par deux bouts. Toute proposition reprise d'une source doit la citer, et
%     toute proposition doit pouvoir etre redemontree au tableau.
% ====================================================================================

\bigskip

\noindent\textit{\myLocation, \myTime}

\smallskip

\begin{flushright}
    \begin{tabular}{m{5cm}}
        \\ \hline
        \centering\myName \\
    \end{tabular}
\end{flushright}
